% poster2_p1.tex / poster2_p2.tex（A0 2 枚版）で共有する設定。poster.tex（1 枚版）の設定を元に、文字を大きくしたもの
% 寸法・文字サイズは A1 のときの値（poster2_a0.tex で 1.414 倍して A0 にする）
\usepackage{xeCJK}
\setCJKmainfont[BoldFont=HaranoAjiGothic-Bold.otf]{HaranoAjiGothic-Regular.otf}
\setCJKsansfont[BoldFont=HaranoAjiGothic-Bold.otf]{HaranoAjiGothic-Regular.otf}
\setmainfont{texgyreheros-regular.otf}[BoldFont=texgyreheros-bold.otf]
\setsansfont{texgyreheros-regular.otf}[BoldFont=texgyreheros-bold.otf]
\xeCJKsetup{CJKspace=false}

\usepackage{amsmath,amssymb,bm}
\usepackage{enumitem}
\usepackage{graphicx}
\usepackage{array}
\usepackage{fontawesome5}
\usepackage{pgfplots}
\pgfplotsset{compat=1.18}
\usetikzlibrary{positioning,arrows.meta,calc}
\graphicspath{{../../resume/}{../../}}

\setlist{nosep,leftmargin=1.0em,itemsep=4pt}
\renewcommand{\labelitemi}{\textcolor{blue}{\textbullet}}
\setlength{\parindent}{0pt}
\setlength{\parskip}{2pt}

% ---------------------------------------------------------------- 色（DESIGN.md とロゴに合わせる）
\definecolor{navy}{HTML}{164A86}     % ロゴの文字色
\definecolor{blue}{HTML}{0093C4}     % blue-deep
\definecolor{skytint}{HTML}{E3F7FF}
\definecolor{yellow}{HTML}{FFDB4F}
\definecolor{yellowdeep}{HTML}{C99A00}
\definecolor{yellowtint}{HTML}{FFF6DA}
\definecolor{ink}{HTML}{232323}
\definecolor{ink2}{HTML}{5E584C}
\definecolor{paper}{HTML}{F1EFE8}    % ロゴの背景色
\definecolor{card}{HTML}{F6F3EA}
\definecolor{good}{HTML}{1C7A45}
\definecolor{goodtint}{HTML}{E3F4E9}
\definecolor{orange}{HTML}{E2612B}

\definecolorstyle{gemmba}{
  \colorlet{colorOne}{navy}\colorlet{colorTwo}{blue}\colorlet{colorThree}{yellow}
}{
  \colorlet{backgroundcolor}{paper}
  \colorlet{framecolor}{paper}
  \colorlet{titlefgcolor}{navy}
  \colorlet{titlebgcolor}{paper}
  \colorlet{blocktitlebgcolor}{white}
  \colorlet{blocktitlefgcolor}{navy}
  \colorlet{blockbodybgcolor}{white}
  \colorlet{blockbodyfgcolor}{ink}
  \colorlet{innerblocktitlebgcolor}{white}
  \colorlet{innerblocktitlefgcolor}{navy}
  \colorlet{innerblockbodybgcolor}{white}
  \colorlet{innerblockbodyfgcolor}{ink}
  \colorlet{notefgcolor}{ink}
  \colorlet{notebgcolor}{skytint}
  \colorlet{noteframecolor}{blue}
}
\usecolorstyle{gemmba}
\usetitlestyle{Filled}
\usebackgroundstyle{Default}

% ブロック：白いカード。見出しは中央、黄色の下線
\defineblockstyle{gemblock}{
  titlewidthscale=1, bodywidthscale=1, titlecenter,
  titleoffsetx=0pt, titleoffsety=0pt, bodyoffsetx=0mm, bodyoffsety=0mm,
  bodyverticalshift=0mm, roundedcorners=10, linewidth=0pt,
  titleinnersep=3mm, bodyinnersep=6mm
}{
  \ifBlockHasTitle
    \draw[draw=none, fill=white, rounded corners=\blockroundedcorners]
      (blockbody.south west) rectangle (blocktitle.north east);
  \else
    \draw[draw=none, fill=white, rounded corners=\blockroundedcorners]
      (blockbody.south west) rectangle (blockbody.north east);
  \fi
}
\useblockstyle{gemblock}


% ---------------------------------------------------------------- タイトル（2 枚とも同じ帯。右上にページ番号）
% \posttitle{ページの主題}{n/2}
\newcommand{\posttitle}[2]{%
  \settitle{%
    \begin{minipage}[c]{0.36\linewidth}
      \includegraphics[width=0.95\linewidth]{logo.png}
    \end{minipage}\hfill
    \begin{minipage}[c]{0.62\linewidth}
      \raggedright\color{navy}
      {\fontsize{66pt}{70pt}\selectfont\bfseries 中小企業向けのDX推進システム}\par\vspace{10pt}
      {\fontsize{36pt}{40pt}\selectfont\bfseries\color{blue}#1}\quad
      {\fontsize{26pt}{30pt}\selectfont\bfseries\color{white}\colorbox{navy}{\,#2\,}}\par\vspace{10pt}
      {\fontsize{24pt}{28pt}\selectfont\color{ink2}明石工業高等専門学校\quad 小笠原\,悠太\quad 増永\,惺真\qquad 指導教員：野村\,隼人}
    \end{minipage}}}

% ---------------------------------------------------------------- 部品
% ブロック見出し：文字幅に合わせた黄色の下線
\newcommand{\bt}[1]{\tikz[baseline=(t.base)]{%
  \node[inner sep=0, font=\fontsize{48pt}{54pt}\selectfont\bfseries, text=navy] (t) {#1};
  \fill[yellow] ([xshift=-6mm, yshift=-16pt]t.south west) rectangle ([xshift=6mm, yshift=-6pt]t.south east);}}
\newcommand{\sm}{\fontsize{20pt}{25pt}\selectfont}
\newcommand{\capt}[1]{{\centering\fontsize{20pt}{25pt}\selectfont\color{ink2}#1\par}}
\newcommand{\head}[1]{{\fontsize{32pt}{38pt}\selectfont\bfseries\color{navy}#1}\par\vspace{10pt}}
\newcommand{\hl}[1]{\textcolor{blue}{\textbf{#1}}}
\newcommand{\bl}{\fontsize{24pt}{31pt}\selectfont}   % 本文・箇条書き
\newcommand{\tb}{\fontsize{25pt}{30pt}\selectfont}   % 表
\renewcommand{\arraystretch}{1.45}
\setlist{nosep,leftmargin=1.0em,itemsep=7pt}

% 写真：img/<ファイル名> があれば表示、無ければ「ここに写真」の枠を描く
%   \photo{ファイル名}{幅}{高さ}{何の写真か}
\newcommand{\photo}[4]{%
  \IfFileExists{img/#1}{%
    \tikz{\node[inner sep=0] {\includegraphics[width=#2,height=#3,keepaspectratio]{img/#1}};}}{%
    \tikz{\draw[dashed, line width=2pt, draw=ink2!55, fill=card, rounded corners=8pt]
            (0,0) rectangle (#2,#3);
          \node[align=center, text=ink2, font=\fontsize{20pt}{25pt}\selectfont]
            at ($(0,0)!0.5!(#2,#3)$)
            {{\fontsize{36pt}{36pt}\selectfont\faCamera}\\[4pt]#4\\{\fontsize{15pt}{18pt}\selectfont img/\detokenize{#1}}};}}}
% 管理画面のスクリーンショット（細い枠つき）
\newcommand{\shot}[2]{\fcolorbox{ink2!40}{white}{\includegraphics[width=#1]{img/#2}}}

% メリット：短い見出し＋補足 1 行。項目内は詰め、項目の間を空ける
\newcommand{\merit}[2]{% 見出し, 補足
  {\bfseries\fontsize{30pt}{36pt}\selectfont #1\par}
  \vspace{3pt}
  {\fontsize{23pt}{28pt}\selectfont\color{ink2}#2\par}
  \vspace{14pt}}

% 囲み（淡い色の地のカード）\tint{地の色}{中身}
\newcommand{\tint}[2]{\coloredbox[bgcolor=#1, fgcolor=ink, framecolor=#1, roundedcorners=8, linewidth=0pt]{#2}}

\tikzset{
  box/.style={draw=navy, line width=3pt, rounded corners=8pt, fill=white,
    align=flush center, inner sep=9pt, text=ink},
  aibox/.style={draw=blue, line width=3.5pt, rounded corners=8pt, fill=skytint,
    align=flush center, inner sep=9pt, text=ink},
  ybox/.style={draw=yellowdeep, line width=3.5pt, rounded corners=8pt, fill=yellowtint,
    align=flush center, inner sep=9pt, text=ink},
  gbox/.style={draw=good, line width=3.5pt, rounded corners=8pt, fill=goodtint,
    align=flush center, inner sep=9pt, text=ink},
  obox/.style={draw=orange, line width=3.5pt, rounded corners=8pt, fill=orange!8,
    align=flush center, inner sep=9pt, text=ink},
  arr/.style={-{Stealth[length=18pt,width=17pt]}, line width=4pt, color=navy!70},
  dbl/.style={{Stealth[length=15pt,width=14pt]}-{Stealth[length=15pt,width=14pt]},
    line width=3.5pt, color=navy!70},
  lab/.style={font=\fontsize{20pt}{24pt}\selectfont, text=ink, align=left},
  elab/.style={font=\fontsize{19pt}{22pt}\selectfont, text=ink2},
  slab/.style={font=\fontsize{17pt}{20pt}\selectfont, text=ink2},   % 矢印が短いところのラベル
}

\pgfplotsset{
  gem/.style={
    axis lines=left, axis line style={ink2, line width=1.2pt},
    tick style={ink2, line width=1.2pt},
    tick label style={font=\fontsize{20pt}{22pt}\selectfont, color=ink},
    label style={font=\fontsize{20pt}{23pt}\selectfont, color=ink2},
    every axis plot/.append style={line width=3.5pt},
  },
}

\colorlet{orangetint}{orange!8}
